% ============================================================
\chapter{Soluções dos exercícios}
\label{ap:solucoes}
% ============================================================

As soluções abaixo são \emph{uma} forma de resolver cada exercício; outras soluções corretas são possíveis e
devem ser valorizadas. Todas as saídas mostradas foram obtidas executando os códigos.

\begin{entrada}
import math
import numpy as np
import sympy as sp
import matplotlib.pyplot as plt
\end{entrada}

% ------------------------------------------------------------
\section*{Aula 1}
% ------------------------------------------------------------

\paragraph{Exercício \ref{ex:1-temp}.}
% entradas: 36.5
\begin{entrada}
C = float(input("Temperatura em °C: "))
F = 9/5*C + 32
print(f"{C} °C = {F:.1f} °F")
\end{entrada}
\begin{saida}
Temperatura em °C: 36.5
36.5 °C = 97.7 °F
\end{saida}

\paragraph{Exercício \ref{ex:1-bissexto}.}
\begin{entrada}
for ano in [1900, 2000, 2024, 2025]:
    bissexto = (ano % 4 == 0 and ano % 100 != 0) or ano % 400 == 0
    print(ano, "é bissexto" if bissexto else "não é bissexto")
\end{entrada}
\begin{saida}
1900 não é bissexto
2000 é bissexto
2024 é bissexto
2025 não é bissexto
\end{saida}

\paragraph{Exercício \ref{ex:1-somas}.}
\begin{entrada}
s = 0
for k in range(1, 100, 2):
    s += k
print("(a)", s)

s = 0
for k in range(1, 21):
    s += 1/k**2
print("(b)", s, "  pi^2/6 =", math.pi**2/6)

n = 10
s = 0
for k in range(1, n + 1):
    s += k**3
print("(c)", s, (n*(n + 1)//2)**2)
\end{entrada}
\begin{saida}
(a) 2500
(b) 1.5961632439130233   pi^2/6 = 1.6449340668482264
(c) 3025 3025
\end{saida}

\paragraph{Exercício \ref{ex:1-5000}.}
\begin{entrada}
soma, k = 0, 0
while soma <= 5000:
    k += 1
    soma += k
print("último número somado:", k, "  soma:", soma)
\end{entrada}
\begin{saida}
último número somado: 100   soma: 5050
\end{saida}

\paragraph{Exercício \ref{ex:1-pa}.}
% entradas: 2 | 3 | 6
\begin{entrada}
a1 = float(input("a1 = "))
r = float(input("razão = "))
n = int(input("n = "))
soma = 0
for k in range(n):
    termo = a1 + k*r
    print(termo, end=" ")
    soma += termo
print("\nsoma =", soma, "  fórmula =", n*(a1 + termo)/2)
\end{entrada}
\begin{saida}
a1 = 2
razão = 3
n = 6
2.0 5.0 8.0 11.0 14.0 17.0
soma = 57.0   fórmula = 57.0
\end{saida}

\paragraph{Exercício \ref{ex:1-collatz}.}
\begin{entrada}
for n in [6, 27]:
    passos = 0
    seq = [n]
    while n != 1:
        n = n // 2 if n % 2 == 0 else 3*n + 1
        seq.append(n)
        passos += 1
    print(passos, "passos; maior valor:", max(seq), "; início:", seq[:8])
\end{entrada}
\begin{saida}
8 passos; maior valor: 16 ; início: [6, 3, 10, 5, 16, 8, 4, 2]
111 passos; maior valor: 9232 ; início: [27, 82, 41, 124, 62, 31, 94, 47]
\end{saida}

\paragraph{Exercício \ref{ex:1-divisores}.}
% entradas: 28
\begin{entrada}
n = int(input("n = "))
qtd = 0
for d in range(1, n + 1):
    if n % d == 0:
        print(d, end=" ")
        qtd += 1
print()
print(n, "é primo" if qtd == 2 else "não é primo")
\end{entrada}
\begin{saida}
n = 28
1 2 4 7 14 28
28 não é primo
\end{saida}

% ------------------------------------------------------------
\section*{Aula 2}
% ------------------------------------------------------------

\paragraph{Exercício \ref{ex:2-listas}.}
\begin{entrada}
L = [3, 8, -2, 7, 0, 5, -6, 1]
print(L[2], L[-1])
print([v for v in L if v > 0])
print([v**2 for v in L])
print(sum(L)/len(L), sorted(L, reverse=True))
\end{entrada}
\begin{saida}
-2 1
[3, 8, 7, 5, 1]
[9, 64, 4, 49, 0, 25, 36, 1]
2.0 [8, 7, 5, 3, 1, 0, -2, -6]
\end{saida}

\paragraph{Exercício \ref{ex:2-quadrados}.}
\begin{entrada}
def soma_quadrados(n):
    return sum([k**2 for k in range(1, n + 1)])

print(all(soma_quadrados(n) == n*(n + 1)*(2*n + 1)//6 for n in range(1, 21)))
\end{entrada}
\begin{saida}
True
\end{saida}

\paragraph{Exercício \ref{ex:2-estat}.}
\begin{entrada}
def media(L):
    return sum(L) / len(L)

def desvio(L):
    m = media(L)
    return math.sqrt(sum([(v - m)**2 for v in L]) / len(L))

notas = [9.5, 7.5, 10, 8.0, 6.5]
print(media(notas), desvio(notas))
\end{entrada}
\begin{saida}
8.3 1.2884098726725126
\end{saida}

\paragraph{Exercício \ref{ex:2-gemeos}.}
\begin{entrada}
def eh_primo(n):
    if n < 2:
        return False
    d = 2
    while d*d <= n:
        if n % d == 0:
            return False
        d += 1
    return True

print([(p, p + 2) for p in range(2, 100) if eh_primo(p) and eh_primo(p + 2)])
\end{entrada}
\begin{saida}
[(3, 5), (5, 7), (11, 13), (17, 19), (29, 31), (41, 43), (59, 61), (71, 73)]
\end{saida}

\paragraph{Exercício \ref{ex:2-mdc}.}
\begin{entrada}
def mdc(a, b):
    while b != 0:
        a, b = b, a % b
    return a

def mmc(a, b):
    return a*b // mdc(a, b)

m = 1
for k in range(2, 11):
    m = mmc(m, k)
print(mdc(252, 198), mmc(4, 6), m)
\end{entrada}
\begin{saida}
18 12 2520
\end{saida}

\paragraph{Exercício \ref{ex:2-poli}.}
\begin{entrada}
def avalia(coef, x):
    return sum([a * x**k for k, a in enumerate(coef)])

def deriva(coef):
    return [k * coef[k] for k in range(1, len(coef))]

p = [1, -3, 0, 2]          # 1 - 3x + 2x^3
print(avalia(p, 2), deriva(p), avalia(deriva(p), 2))
\end{entrada}
\begin{saida}
11 [-3, 0, 6] 21
\end{saida}

% ------------------------------------------------------------
\section*{Aula 3}
% ------------------------------------------------------------

\paragraph{Exercício \ref{ex:3-dist}.}
\begin{entrada}
def distancia(P, Q):
    return math.hypot(Q[0] - P[0], Q[1] - P[1])

A, B, C = (0, 0), (4, 0), (1, 3)
print(distancia(A, B) + distancia(B, C) + distancia(C, A))
\end{entrada}
\begin{saida}
11.404918347287664
\end{saida}

\paragraph{Exercício \ref{ex:3-relogio}.}
\begin{entrada}
def angulo(h, m):
    dif = abs((30*h + 0.5*m) - 6*m) % 360
    return min(dif, 360 - dif)

print(angulo(3, 0), angulo(9, 45), angulo(12, 30))
\end{entrada}
\begin{saida}
90.0 22.5 165.0
\end{saida}

\paragraph{Exercício \ref{ex:3-series}.}
\begin{entrada}
for n in [5, 10, 15]:
    s = sum([1/math.factorial(k) for k in range(n + 1)])
    print(n, s, abs(s - math.e))
for n in [10, 1000, 100000]:
    s = 4*sum([(-1)**k/(2*k + 1) for k in range(n + 1)])
    print(n, s, abs(s - math.pi))
\end{entrada}
\begin{saida}
5 2.716666666666667 0.0016151617923783057
10 2.7182818011463845 2.7312660577649694e-08
15 2.7182818284589945 5.062616992290714e-14
10 3.232315809405593 0.09072315581579993
1000 3.1425916543395434 0.000999000749750234
100000 3.141602653489794 9.999900000945416e-06
\end{saida}

A série de $e$ converge muito mais rápido: o erro cai como $\frac{1}{(n+1)!}$, enquanto na de Leibniz cai como
$\frac1n$.

\paragraph{Exercício \ref{ex:3-seno}.}
\begin{entrada}
def seno(x, n):
    return sum([(-1)**k * x**(2*k + 1) / math.factorial(2*k + 1) for k in range(n + 1)])

for x in [1, 10]:
    print(x, seno(x, 5), seno(x, 20), math.sin(x))
\end{entrada}
\begin{saida}
1 0.841470984648068 0.8414709848078965 0.8414709848078965
10 -1056.9392336059002 -0.5440211107317661 -0.5440211108893698
\end{saida}

\paragraph{Exercício \ref{ex:3-cubica}.}
\begin{entrada}
def raiz_cubica(x, tol=1e-12):
    y = x if abs(x) > 1 else 1.0
    while True:
        novo = (2*y + x/y**2) / 3
        if abs(novo - y) < tol:
            return novo
        y = novo

print(raiz_cubica(27), raiz_cubica(2), math.cbrt(2), raiz_cubica(-8))
\end{entrada}
\begin{saida}
3.0 1.2599210498948732 1.2599210498948734 -2.0
\end{saida}

\paragraph{Exercício \ref{ex:3-bissec}.}
\begin{entrada}
def bisseccao(f, a, b, tol=1e-10):
    while (b - a)/2 > tol:
        m = (a + b)/2
        if f(a)*f(m) <= 0:
            b = m
        else:
            a = m
    return (a + b)/2

print(bisseccao(lambda x: x**2 - 2, 1, 2), math.sqrt(2))
print(bisseccao(lambda x: math.exp(-x) - x, 0, 1))
print(bisseccao(lambda x: x**7 + x - 1, 0, 1))
\end{entrada}
\begin{saida}
1.4142135623260401 1.4142135623730951
0.5671432903618552
0.7965443541179411
\end{saida}

% ------------------------------------------------------------
\section*{Aula 4}
% ------------------------------------------------------------

\paragraph{Exercício \ref{ex:4-criar}.}
\begin{entrada}
print(np.arange(1, 40, 2))
print(np.linspace(-1, 1, 11))
print(2**np.arange(11))
print(1/np.arange(1, 11))
\end{entrada}
\begin{saida}
[ 1  3  5  7  9 11 13 15 17 19 21 23 25 27 29 31 33 35 37 39]
[-1.  -0.8 -0.6 -0.4 -0.2  0.   0.2  0.4  0.6  0.8  1. ]
[   1    2    4    8   16   32   64  128  256  512 1024]
[1.         0.5        0.33333333 0.25       0.2        0.16666667
 0.14285714 0.125      0.11111111 0.1       ]
\end{saida}

\paragraph{Exercício \ref{ex:4-tabela}.}
\begin{entrada}
x = np.linspace(0, np.pi, 7)
for xi, s, c in zip(x, np.sin(x), np.cos(x)):
    print(f"{xi:6.3f} {s:8.4f} {c:8.4f} {s**2 + c**2:6.2f}")
\end{entrada}
\begin{saida}
 0.000   0.0000   1.0000   1.00
 0.524   0.5000   0.8660   1.00
 1.047   0.8660   0.5000   1.00
 1.571   1.0000   0.0000   1.00
 2.094   0.8660  -0.5000   1.00
 2.618   0.5000  -0.8660   1.00
 3.142   0.0000  -1.0000   1.00
\end{saida}

\paragraph{Exercício \ref{ex:4-estat}.}
\begin{entrada}
h = np.array([1.72, 1.65, 1.80, 1.58, 1.75, 1.69, 1.83, 1.61])
print(h.mean(), h.std(), np.sum(h > 1.70), h[h > 1.70].mean())
\end{entrada}
\begin{saida}
1.70375 0.08275528684017716 4 1.775
\end{saida}

\paragraph{Exercício \ref{ex:4-harmonica}.}
\begin{entrada}
k = np.arange(1, 100001)
H = np.cumsum(1/k)
n = np.argmax(H > 10) + 1          # primeiro índice em que passa de 10
print(n, H[n-1])
N = 10**6
print(np.sum(1/np.arange(1, N + 1)) - np.log(N))
\end{entrada}
\begin{saida}
12367 10.000043008275778
0.5772161649014524
\end{saida}

\paragraph{Exercício \ref{ex:4-simpson}.}
\begin{entrada}
def simpson(f, a, b, n):
    x = np.linspace(a, b, n + 1)
    y = f(x)
    h = (b - a)/n
    return h/3 * (y[0] + 4*np.sum(y[1:-1:2]) + 2*np.sum(y[2:-1:2]) + y[-1])

def trapezios(f, a, b, n):
    x = np.linspace(a, b, n + 1)
    y = f(x)
    return (b - a)/n * (np.sum(y) - (y[0] + y[-1])/2)

g = lambda x: 4/(1 + x**2)
for n in [10, 20, 40]:
    print(n, abs(trapezios(g, 0, 1, n) - np.pi), abs(simpson(g, 0, 1, n) - np.pi))
\end{entrada}
\begin{saida}
10 0.0016666646826344333 3.9650577932093256e-08
20 0.00041666663566441997 6.200080449048073e-10
40 0.00010416666618251469 9.688250202088966e-12
\end{saida}

\paragraph{Exercício \ref{ex:4-mc}.}
\begin{entrada}
rng = np.random.default_rng(0)
u = rng.random(100000)
print(np.mean(np.exp(-u**2)), trapezios(lambda x: np.exp(-x**2), 0, 1, 100))
\end{entrada}
\begin{saida}
0.7471060128813966 0.7468180014679697
\end{saida}

% ------------------------------------------------------------
\section*{Aula 5}
% ------------------------------------------------------------

\paragraph{Exercício \ref{ex:5-ops}.}
\begin{entrada}
A = np.array([[1, 2, 0], [0, 1, 3], [4, 0, 1]])
B = np.array([[2, 0, 1], [1, 1, 0], [0, 3, 1]])
print(A + B)
print(A @ B)
print(B @ A)
print(A.T @ B)
print(np.linalg.det(A @ B), np.linalg.det(A) * np.linalg.det(B))
print(np.array_equal((A @ B).T, B.T @ A.T))
\end{entrada}
\begin{saida}
[[3 2 1]
 [1 2 3]
 [4 3 2]]
[[ 4  2  1]
 [ 1 10  3]
 [ 8  3  5]]
[[ 6  4  1]
 [ 1  3  3]
 [ 4  3 10]]
[[ 2 12  5]
 [ 5  1  2]
 [ 3  6  1]]
125.00000000000004 124.99999999999996
True
\end{saida}

\paragraph{Exercício \ref{ex:5-sistema}.} (Usa a função \py{gauss} da Aula 5.)
\begin{entrada}
def gauss(A, b):
    M = np.column_stack([A, b]).astype(float)
    n = len(b)
    for i in range(n - 1):
        p = i + np.argmax(np.abs(M[i:, i]))
        M[[i, p]] = M[[p, i]]
        for j in range(i + 1, n):
            M[j] = M[j] - M[j, i]/M[i, i] * M[i]
    x = np.zeros(n)
    for i in range(n - 1, -1, -1):
        x[i] = (M[i, -1] - M[i, i+1:n] @ x[i+1:]) / M[i, i]
    return x

A = np.array([[2, 1, -1], [-3, -1, 2], [-2, 1, 2]])
b = np.array([8, -11, -3])
print(np.linalg.solve(A, b), gauss(A, b))
\end{entrada}
\begin{saida}
[ 2.  3. -1.] [ 2.  3. -1.]
\end{saida}

\paragraph{Exercício \ref{ex:5-hilbert}.}
\begin{entrada}
for n in [4, 8, 12]:
    i = np.arange(1, n + 1).reshape(n, 1)
    H = 1/(i + i.T - 1)
    x = np.ones(n)
    y = np.linalg.solve(H, H @ x)
    print(n, f"erro máximo = {np.max(np.abs(y - x)):.1e}", f"cond = {np.linalg.cond(H):.1e}")
\end{entrada}
\begin{saida}
4 erro máximo = 5.4e-14 cond = 1.6e+04
8 erro máximo = 1.3e-07 cond = 1.5e+10
12 erro máximo = 2.8e-01 cond = 1.6e+16
\end{saida}

A matriz de Hilbert é mal condicionada: pequenos erros de arredondamento são ampliados por um fator da ordem do
número de condição, e para $n=12$ a solução perde quase todos os algarismos corretos.

\paragraph{Exercício \ref{ex:5-autov}.}
\begin{entrada}
A = np.array([[4, 1], [2, 3]])
lam, V = np.linalg.eig(A)
lam = lam.real
V = V.real
for k in range(2):
    print(lam[k], np.allclose(A @ V[:, k], lam[k] * V[:, k]))
print(lam.sum(), np.trace(A), lam.prod(), np.linalg.det(A))
\end{entrada}
\begin{saida}
5.0 True
2.0 True
7.0 7 10.0 10.000000000000002
\end{saida}

\paragraph{Exercício \ref{ex:5-poli}.}
\begin{entrada}
for p in [[1, 0, -5, 0, 4], [1, 0, -2, -5]]:
    r = np.roots(p)
    print(r, np.round(np.abs(np.polyval(p, r)), 10))
\end{entrada}
\begin{saida}
[-2.  2. -1.  1.] [0. 0. 0. 0.]
[ 2.09455148+0.j         -1.04727574+1.13593989j -1.04727574-1.13593989j] [0. 0. 0.]
\end{saida}

\paragraph{Exercício \ref{ex:5-ajuste}.}
\begin{entrada}
t = np.array([0, 0.5, 1, 1.5, 2, 2.5])
h = np.array([1.0, 7.9, 12.4, 14.6, 14.3, 11.7])
a, b, c = np.polyfit(t, h, 2)
print(f"h(t) = {a:.3f} t^2 + {b:.3f} t + {c:.3f}")
print(f"g ~ {-2*a:.2f} m/s^2,  altura máxima em t = {-b/(2*a):.2f} s")
\end{entrada}
\begin{saida}
h(t) = -4.764 t^2 + 16.191 t + 0.996
g ~ 9.53 m/s^2,  altura máxima em t = 1.70 s
\end{saida}

% ------------------------------------------------------------
\section*{Aula 6}
% ------------------------------------------------------------

\paragraph{Exercício \ref{ex:6-derivada}.}
% figura: sol-6-1
\begin{entrada}
x = np.linspace(-2.5, 2.5, 300)
plt.plot(x, x**3 - 3*x, label="$f(x)=x^3-3x$")
plt.plot(x, 3*x**2 - 3, "--", label="$f'(x)=3x^2-3$")
plt.plot([-1, 1], [2, -2], "ro", label="pontos críticos")
plt.axhline(0, color="k", lw=0.8); plt.axvline(0, color="k", lw=0.8)
plt.grid(True); plt.legend()
plt.show()
\end{entrada}
\saidafig[0.5]{sol-6-1}

\paragraph{Exercício \ref{ex:6-explog}.}
\begin{entrada}
fig, ax = plt.subplots(1, 2, figsize=(10, 4))
x = np.linspace(-2, 2, 200)
for base, nome in [(2, "$2^x$"), (np.e, "$e^x$"), (10, "$10^x$")]:
    ax[0].plot(x, base**x, label=nome)
ax[0].set_ylim(0, 8); ax[0].legend()
x = np.linspace(0.05, 10, 300)
ax[1].plot(x, np.log2(x), label=r"$\log_2 x$")
ax[1].plot(x, np.log(x), label=r"$\ln x$")
ax[1].plot(x, np.log10(x), label=r"$\log_{10} x$")
ax[1].legend()
plt.show()
\end{entrada}

\paragraph{Exercício \ref{ex:6-parabolas}.} O vértice de $y=x^2+bx+1$ é $\left(-\frac b2, 1-\frac{b^2}{4}\right)$;
escrevendo $x_v=-\frac b2$, temos $y_v=1-x_v^2$: os vértices estão sobre a parábola $y=1-x^2$.
\begin{entrada}
x = np.linspace(-4, 4, 300)
for b in range(-3, 4):
    plt.plot(x, x**2 + b*x + 1, lw=1)
    plt.plot(-b/2, 1 - b**2/4, "ko")
plt.plot(x, 1 - x**2, "r--", label="$y = 1 - x^2$")
plt.ylim(-3, 6); plt.legend()
plt.show()
\end{entrada}

\paragraph{Exercício \ref{ex:6-param}.}
% figura: sol-6-4
\begin{entrada}
t = np.linspace(0, 2*np.pi, 400)
fig = plt.figure(figsize=(11, 3.5))
ax1 = fig.add_subplot(1, 3, 1)
ax1.plot(np.cos(t)**3, np.sin(t)**3); ax1.set_aspect("equal"); ax1.set_title("astroide")
ax2 = fig.add_subplot(1, 3, 2)
X, Y = np.meshgrid(np.linspace(-1.2, 1.2, 400), np.linspace(-0.6, 0.6, 300))
ax2.contour(X, Y, (X**2 + Y**2)**2 - (X**2 - Y**2), levels=[0])
ax2.set_aspect("equal"); ax2.set_title("lemniscata")
ax3 = fig.add_subplot(1, 3, 3, projection="polar")
ax3.plot(t, np.sin(4*t)); ax3.set_title(r"$r = \sin 4\theta$")
plt.show()
\end{entrada}
\saidafig[0.95]{sol-6-4}

\paragraph{Exercício \ref{ex:6-sup}.} O máximo fica perto de $\left(\frac{1}{\sqrt2},0\right)$ e o mínimo perto de
$\left(-\frac{1}{\sqrt2},0\right)$ (confirme na Aula 8 igualando as derivadas parciais a zero).
\begin{entrada}
x = np.linspace(-2, 2, 100)
X, Y = np.meshgrid(x, x)
Z = X*np.exp(-X**2 - Y**2)
fig = plt.figure(figsize=(10, 4))
ax1 = fig.add_subplot(1, 2, 1, projection="3d")
ax1.plot_surface(X, Y, Z, cmap="coolwarm")
ax2 = fig.add_subplot(1, 2, 2)
c = ax2.contour(X, Y, Z, levels=15, cmap="coolwarm")
ax2.clabel(c, fontsize=7); ax2.set_aspect("equal")
plt.show()
\end{entrada}

\paragraph{Exercício \ref{ex:6-secantes}.}
\begin{entrada}
f = np.sqrt
x = np.linspace(0, 4, 200)
plt.plot(x, f(x), "k", lw=2, label=r"$\sqrt{x}$")
for hh in [2, 1, 0.5, 0.1]:
    m = (f(1 + hh) - f(1)) / hh
    plt.plot(x, f(1) + m*(x - 1), lw=1, label=f"secante h = {hh}")
plt.plot(x, 1 + (x - 1)/2, "r--", lw=2, label="tangente")
plt.ylim(0, 2.5); plt.legend(fontsize=8)
plt.show()
\end{entrada}

% ------------------------------------------------------------
\section*{Aula 7}
% ------------------------------------------------------------

\begin{entrada}
x, y, a, b, c = sp.symbols('x y a b c')
\end{entrada}

\paragraph{Exercício \ref{ex:7-exato}.}
\begin{entrada}
for e in [sp.sqrt(12) + sp.sqrt(27), sp.Rational(1, 3) + sp.Rational(2, 7), sp.exp(sp.pi*sp.sqrt(163))]:
    print(e, "=", sp.N(e, 30))
\end{entrada}
\begin{saida}
5*sqrt(3) = 8.66025403784438646763723170753
13/21 = 0.619047619047619047619047619048
exp(sqrt(163)*pi) = 262537412640768743.999999999999
\end{saida}

\paragraph{Exercício \ref{ex:7-manip}.}
\begin{entrada}
print(sp.expand((x - 2*y)**6))
print(sp.factor(x**6 - 1), "  |  ", sp.factor(x**4 - 5*x**2 + 4))
print(sp.cancel((x**3 - 8)/(x**2 - 4)))     # simplify não reduz esta fração; cancel sim
print(sp.apart((x + 1)/(x**2*(x - 1))))
\end{entrada}
\begin{saida}
x**6 - 12*x**5*y + 60*x**4*y**2 - 160*x**3*y**3 + 240*x**2*y**4 - 192*x*y**5 + 64*y**6
(x - 1)*(x + 1)*(x**2 - x + 1)*(x**2 + x + 1)   |   (x - 2)*(x - 1)*(x + 1)*(x + 2)
(x**2 + 2*x + 4)/(x + 2)
2/(x - 1) - 2/x - 1/x**2
\end{saida}

\paragraph{Exercício \ref{ex:7-eqs}.} Em (c), elevar ao quadrado dá $x^2-7x+6=0$, com raízes 1 e 6; mas $x=1$ não
satisfaz a equação original (é uma raiz ``estranha''). O \py{solve} já descarta essa raiz.
\begin{entrada}
print(sp.solve(x**3 - 6*x**2 + 11*x - 6, x))
print(sp.solve(2**x - 32, x))
print(sp.solve(sp.Eq(sp.sqrt(x + 3), x - 3), x))
print(sp.solveset(x**4 + x**2 - 12, x, sp.S.Reals), sp.solveset(x**4 + x**2 - 12, x, sp.S.Complexes))
print(sp.solveset(sp.Abs(2*x - 1) < 3, x, sp.S.Reals))
\end{entrada}
\begin{saida}
[1, 2, 3]
[5]
[6]
{-sqrt(3), sqrt(3)} {-sqrt(3), sqrt(3), -2*I, 2*I}
Interval.open(-1, 2)
\end{saida}

\paragraph{Exercício \ref{ex:7-sist}.}
\begin{entrada}
z = sp.symbols('z')
print(sp.linsolve([x + 2*y - z - 4, 2*x - y + z - 1, x + y + z - 6], [x, y, z]))
print(sp.solve([x**2 + y**2 - 25, x - y - 1], [x, y], dict=True))
p = a*x**2 + b*x + c
print(sp.solve([p.subs(x, 1) - 2, p.subs(x, 2) - 3, p.subs(x, 3) - 6], [a, b, c]))
\end{entrada}
\begin{saida}
{(5/7, 20/7, 17/7)}
[{x: -3, y: -4}, {x: 4, y: 3}]
{a: 1, b: -2, c: 3}
\end{saida}

\paragraph{Exercício \ref{ex:7-nsolve}.} O gráfico (\py{sp.plot(sp.exp(-x) - x**2 + 2, (x, -3, 3))}) mostra duas
raízes, perto de $-1$ e de $1{,}5$.
\begin{entrada}
f = sp.exp(-x) - x**2 + 2
print(sp.nsolve(f, x, -1), sp.nsolve(f, x, 1.5))
\end{entrada}
\begin{saida}
1.49164419466824 1.49164419466824
\end{saida}

\paragraph{Exercício \ref{ex:7-somas}.}
\begin{entrada}
k, n = sp.symbols('k n', integer=True, positive=True)
print(sp.factor(sp.summation(2*k - 1, (k, 1, n))), "|", sp.factor(sp.summation(k*(k + 1), (k, 1, n))),
      "|", sp.summation(2**k, (k, 0, n)))
print(sp.summation(1/(k*(k + 1)), (k, 1, sp.oo)), sp.summation(1/sp.factorial(k), (k, 0, sp.oo)))
\end{entrada}
\begin{saida}
n**2 | n*(n + 1)*(n + 2)/3 | 2**(n + 1) - 1
1 E
\end{saida}

% ------------------------------------------------------------
\section*{Aula 8}
% ------------------------------------------------------------

\paragraph{Exercício \ref{ex:8-limites}.}
\begin{entrada}
print(sp.limit((sp.tan(x) - x)/x**3, x, 0, '+-'),
      sp.limit((1 + 2/x)**x, x, sp.oo),
      sp.limit(x*sp.log(x), x, 0, '+'),
      sp.limit(sp.sqrt(x**2 + x) - x, x, sp.oo),
      sp.limit((sp.sqrt(x + 1) - 2)/(x - 3), x, 3, '+-'))
\end{entrada}
\begin{saida}
1/3 exp(2) 0 1/2 1/4
\end{saida}

\paragraph{Exercício \ref{ex:8-cont}.}
\begin{entrada}
kk = sp.symbols('k')
esq = sp.limit(kk*x**2, x, 2, '-')
dir_ = sp.limit(2*x + kk, x, 2, '+')
print(esq, dir_, sp.solve(sp.Eq(esq, dir_), kk))
\end{entrada}
\begin{saida}
4*k k + 4 [4/3]
\end{saida}

\paragraph{Exercício \ref{ex:8-derivadas}.}
\begin{entrada}
for f in [x**x, sp.log((1 + x)/(1 - x)), sp.exp(-x**2)*sp.cos(3*x), sp.asin(sp.sqrt(x))]:
    print(sp.simplify(sp.diff(f, x)))
print(sp.diff(sp.exp(-x**2)*sp.cos(3*x), x).subs(x, 0))
\end{entrada}
\begin{saida}
x**x*(log(x) + 1)
-2/(x**2 - 1)
-(2*x*cos(3*x) + 3*sin(3*x))*exp(-x**2)
1/(2*sqrt(x)*sqrt(1 - x))
0
\end{saida}

\paragraph{Exercício \ref{ex:8-tangente}.}
\begin{entrada}
f = x/(1 + x**2)
m = sp.diff(f, x).subs(x, 1)
y0 = f.subs(x, 1)
print("m =", m, "  tangente: y =", y0 + m*(x - 1), "  normal: x = 1" if m == 0 else "")
\end{entrada}
\begin{saida}
m = 0   tangente: y = 1/2   normal: x = 1
\end{saida}

Como $f'(1)=0$, a tangente é horizontal, $y=\frac12$, e a normal é a reta vertical $x=1$.

\paragraph{Exercício \ref{ex:8-estudo}.}
\begin{entrada}
f = x**4 - 4*x**3 + 10
df, d2f = sp.diff(f, x), sp.diff(f, x, 2)
print("críticos:", sp.solve(df, x), "  inflexão (candidatos):", sp.solve(d2f, x))
for c in sp.solve(df, x):
    print(c, f.subs(x, c), d2f.subs(x, c))
print(sp.solveset(df > 0, x, sp.S.Reals))
\end{entrada}
\begin{saida}
críticos: [0, 3]   inflexão (candidatos): [0, 2]
0 10 0
3 -17 36
Interval.open(3, oo)
\end{saida}

$f$ decresce em $(-\infty,3)$ e cresce em $(3,\infty)$; mínimo em $x=3$ ($f=-17$); em $x=0$ há ponto crítico sem
extremo (é ponto de inflexão com tangente horizontal); as inflexões estão em $x=0$ e $x=2$.

\paragraph{Exercício \ref{ex:8-otim}.}
\begin{entrada}
r = sp.symbols('r', positive=True)
A = 2*sp.pi*r**2 + 2*sp.pi*r*1000/(sp.pi*r**2)
rc = sp.solve(sp.diff(A, r), r)[0]
h = 1000/(sp.pi*rc**2)
print(rc, "~", rc.evalf(4), "cm;  h ~", h.evalf(4), "cm;  h/r =", sp.simplify(h/rc))
\end{entrada}
\begin{saida}
5*2**(2/3)/pi**(1/3) ~ 5.419 cm;  h ~ 10.84 cm;  h/r = 2
\end{saida}

A lata ótima tem altura igual ao diâmetro ($h=2r$).

\paragraph{Exercício \ref{ex:8-taylor}.}
\begin{entrada}
P5 = sp.series(sp.sin(x), x, 0, 6).removeO()
aprox = P5.subs(x, sp.Rational(1, 2))
erro = abs(sp.sin(sp.Rational(1, 2)) - aprox).evalf()
print(P5, "|", aprox.evalf(12), "| erro:", erro, "| estimativa:", (sp.Rational(1, 2)**7/sp.factorial(7)).evalf())
\end{entrada}
\begin{saida}
x**5/120 - x**3/6 + x | 0.479427083333 | erro: 1.54472913033306e-6 | estimativa: 1.55009920634921e-6
\end{saida}

% ------------------------------------------------------------
\section*{Aula 9}
% ------------------------------------------------------------

\paragraph{Exercício \ref{ex:9-primitivas}.}
\begin{entrada}
for f in [x**2*sp.exp(-x), (x + 3)/(x**2 + x - 2), sp.sec(x)**2, x*sp.log(x), 1/(1 + 4*x**2)]:
    F = sp.integrate(f, x)
    print(F, "   ok" if sp.simplify(sp.diff(F, x) - f) == 0 else "   conferir")
\end{entrada}
\begin{saida}
(-x**2 - 2*x - 2)*exp(-x)    ok
4*log(x - 1)/3 - log(x + 2)/3    ok
sin(x)/cos(x)    ok
x**2*log(x)/2 - x**2/4    ok
atan(2*x)/2    ok
\end{saida}

\paragraph{Exercício \ref{ex:9-definidas}.}
\begin{entrada}
print(sp.integrate(sp.sin(x)**2, (x, 0, sp.pi/2)), sp.integrate(sp.log(x)/x, (x, 1, sp.E)),
      sp.integrate(x*sp.exp(-x), (x, 0, sp.oo)), sp.integrate(1/sp.sqrt(x), (x, 0, 1)),
      sp.integrate(1/x, (x, 0, 1)))
\end{entrada}
\begin{saida}
pi/4 1/2 1 2 oo
\end{saida}

\paragraph{Exercício \ref{ex:9-area}.} As curvas se cruzam em $x=-2,0,2$; por simetria, a área é o dobro da área
em $[0,2]$.
\begin{entrada}
f, g = 3*x, x**3 - x
print(sp.solve(f - g, x))
area = sp.integrate(g - f, (x, -2, 0)) + sp.integrate(f - g, (x, 0, 2))
print(area, sp.integrate(sp.Abs(f - g), (x, -2, 2)))
\end{entrada}
\begin{saida}
[-2, 0, 2]
8 8
\end{saida}

\paragraph{Exercício \ref{ex:9-riemann}.}
\begin{entrada}
for N in [10, 100, 1000]:
    dx = 4/N
    xe = np.linspace(0, 4 - dx, N)
    somas = [np.sum(np.sqrt(xi))*dx for xi in (xe, xe + dx/2, xe + dx)]
    print(N, ["%.6f" % abs(s - 16/3) for s in somas])
\end{entrada}
\begin{saida}
10 ['0.449259', '0.013737', '0.350741']
100 ['0.041630', '0.000470', '0.038370']
1000 ['0.004052', '0.000015', '0.003948']
\end{saida}

O ponto médio tem o menor erro.

\paragraph{Exercício \ref{ex:9-volume}.}
\begin{entrada}
R, hh = sp.symbols('R h', positive=True)
print(sp.pi*sp.integrate(sp.exp(-x)**2, (x, 0, sp.oo)), sp.pi*sp.integrate((R/hh*x)**2, (x, 0, hh)))
\end{entrada}
\begin{saida}
pi/2 pi*R**2*h/3
\end{saida}

\paragraph{Exercício \ref{ex:9-extra}.}
\begin{entrada}
print(sp.integrate(sp.sin(x), (x, 0, sp.pi))/sp.pi, sp.integrate(x*y, (y, 0, x), (x, 0, 1)))
\end{entrada}
\begin{saida}
2/pi 1/8
\end{saida}

% ------------------------------------------------------------
\section*{Aula 10}
% ------------------------------------------------------------

\paragraph{Exercício \ref{ex:10-mat}.}
\begin{entrada}
A = sp.Matrix([[2, -1, 0], [-1, 2, -1], [0, -1, 2]])
print(A.det(), A.inv())
print(sp.factor(A.charpoly(x).as_expr()), A.eigenvals())
print(A.is_diagonalizable())
\end{entrada}
\begin{saida}
4 Matrix([[3/4, 1/2, 1/4], [1/2, 1, 1/2], [1/4, 1/2, 3/4]])
(x - 2)*(x**2 - 4*x + 2) {2: 1, 2 - sqrt(2): 1, sqrt(2) + 2: 1}
True
\end{saida}

\paragraph{Exercício \ref{ex:10-param}.}
\begin{entrada}
a = sp.symbols('a')
M = sp.Matrix([[1, 1, 1], [1, a, 1], [1, 1, a]])
bv = sp.Matrix([1, 2, 3])
print(sp.factor(M.det()))
print(M.subs(a, 1).rank(), M.subs(a, 1).row_join(bv).rank())
print(sp.simplify(M.LUsolve(bv)).T)
\end{entrada}
\begin{saida}
(a - 1)**2
1 2
Matrix([[(a - 4)/(a - 1), 1/(a - 1), 2/(a - 1)]])
\end{saida}

Solução única para $a\neq1$; para $a=1$ os postos são 1 e 2: sistema impossível.

\paragraph{Exercício \ref{ex:10-edo1}.}
\begin{entrada}
t = sp.symbols('t')
yf = sp.Function('y')
edos = [(sp.Eq(yf(t).diff(t), -2*t*yf(t)), {yf(0): 1}),
        (sp.Eq(yf(t).diff(t) + yf(t), t), None),
        (sp.Eq(yf(t).diff(t), yf(t)/t + t), None)]
for edo, ci in edos:
    sol = sp.dsolve(edo, yf(t), ics=ci)
    print(sol, sp.checkodesol(edo, sol))
\end{entrada}
\begin{saida}
Eq(y(t), exp(-t**2)) (True, 0)
Eq(y(t), C1*exp(-t) + t - 1) (True, 0)
Eq(y(t), t*(C1 + t)) (True, 0)
\end{saida}

\paragraph{Exercício \ref{ex:10-resfriamento}.}
\begin{entrada}
T = sp.Function('T')
k = sp.symbols('k', positive=True)
sol = sp.dsolve(sp.Eq(T(t).diff(t), -k*(T(t) - 25)), T(t), ics={T(0): 90}).rhs
kv = sp.solve(sp.Eq(sol.subs(t, 5), 70), k)[0]
t40 = sp.solve(sp.Eq(sol.subs(k, kv), 40), t)[0]
print(sol, "|  k =", kv, "~", kv.evalf(4), "|  t =", t40.evalf(4), "min")
\end{entrada}
\begin{saida}
25 + 65*exp(-k*t) |  k = log(13**(1/5)*3**(3/5)/3) ~ 0.07355 |  t = 19.94 min
\end{saida}

\paragraph{Exercício \ref{ex:10-edo2}.}
\begin{entrada}
amort = sp.dsolve(yf(t).diff(t, 2) + yf(t).diff(t) + yf(t), yf(t), ics={yf(0): 1, yf(t).diff(t).subs(t, 0): 0})
livre = sp.dsolve(yf(t).diff(t, 2) + yf(t), yf(t), ics={yf(0): 1, yf(t).diff(t).subs(t, 0): 0})
print(amort)
print(livre)
\end{entrada}
\begin{saida}
Eq(y(t), (sqrt(3)*sin(sqrt(3)*t/2)/3 + cos(sqrt(3)*t/2))*exp(-t/2))
Eq(y(t), cos(t))
\end{saida}

O gráfico pode ser feito com \py{sp.plot(amort.rhs, livre.rhs, (t, 0, 15))}: a solução amortecida oscila com
amplitude decrescente, a outra mantém amplitude 1.

\paragraph{Exercício \ref{ex:10-campo}.} A reta $y=t-1$ é solução: $y'=1$ e $t-y=1$.
\begin{entrada}
sol = sp.dsolve(sp.Eq(yf(t).diff(t), t - yf(t)), yf(t))
print(sol)
\end{entrada}
\begin{saida}
Eq(y(t), C1*exp(-t) + t - 1)
\end{saida}
